\section{Methods}

\subsection{Data}

Data were retrieved for 168,347 emergency ischaemic stroke admissions to hospitals in England and Wales for six calendar years, 2016–2021 inclusive, obtained from the national stroke registry SSNAP. Further details may be found in the supplementary material.

\subsection{Machine learning models}

We trained a set of six independent XGBoost models \cite{chen_xgboost_2016} to predict the likelihood of an ischaemic stroke patient attaining any given mRS threshold at discharge. For each mRS threshold, from mRS 0 to 5 (alive, but severely disabled) we refer to being at or lower than that threshold as a \textit{`same or better'} or '\textit{`good outcome'}. The results from the test set for the model fitted on the first k-fold split was used to report the relationships between feature values and their contribution to the predictions.

Code, and full results, for the machine learning work may be found at \url{https://github.com/samuel-book/thrombolysis_clinical_trials_ml_paper}. Further details of method development may be found in the supplementary material.

\subsection{Feature selection and model building}

The complete available data set contained 57 features that described patient demographic and clinical characteristics, the acute stroke pathway, the use/time of thrombolysis, and the stroke team attended. We have preciously identified features that influence the choice to use thrombolysis \cite{pearn_are_2024}. For the current work, we identified the features that influence patient outcome and, after discussions with clinicians, combined the two sets of features (on use of thrombolysis, and outcome after stroke) into a proposed causal model with seven features that would adequately isolate the effect of thrombolysis on outcome.

\subsection{Model accuracy}

For each of the six machine learning models, we report the model accuracy measured using 5-fold cross validation where the data is split into 5 different 80:20 train/test splits with each patient being in one, and only one, test set.

\subsection{SHapley Additive exPlanation (SHAP) values}

SHAP values quantify the contribution of each feature to the model’s prediction of a good discharge outcome \cite{lundberg_unified_2017}. Because SHAP values are expressed on the log-odds scale, they are additive and therefore more suitable than probabilities for assessing the contribution of individual features to the final prediction. Separate SHAP analyses were performed for each of the six machine learning models, corresponding to the six mRS thresholds. A positive SHAP value indicates that a feature increases the likelihood of reaching the mRS threshold (good outcome), whereas a negative SHAP value indicates that a feature decreases it. SHAP values can be examined both locally, at the individual patient level, and globally, across the cohort, to identify general patterns in how patient, pathway, and hospital characteristics are associated with discharge outcomes.

\subsection{Counterfactual treatment:  What if patients had not received thrombolysis?}

After constructing a model (splitting data 80\% training, 20\% test) to isolate the effect of thrombolysis from confounding variables, we used counterfactuals for each patient in the test set who had received thrombolysis, to assess the contribution that thrombolysis made on the likelihood of a good outcome at discharge. The counterfactual result was achieved by running the same model, but changing the patient to mimic not receiving thrombolysis.

The predicted contribution of thrombolysis to the achievement of an excellent outcome (mRS 0–1) was estimated by the difference in SAHP thrombolysis values between when thrombolysis was administered and the counterfactual scenario of not receiving thrombolysis. To model decay of thrombolysis benefit with treatment delay, we fitted a linear regression of estimated thrombolysis benefit (log-odds shift for achieving mRS 0–1) on time from onset to treatment, excluding patients treated after 300 minutes. We defined an excellent outcome as mRS 0–1 to align with the definition used in the clinical-trial meta-analysis by Emberson et al. \cite{emberson_effect_2014}. Analyses were performed in three thrombolysed cohorts: (1) all ischaemic strokes (n = 6,897), (2) severe ischaemic strokes (NIHSS $\geq$ 11; n = 2,897), and (3) mild-to-moderate ischaemic strokes (NIHSS 0–10; n = 4,000).
